\section{The \textsc{predictor4} Classifier}
\label{sec:classifier}

We describe the architecture of \textsc{predictor4}, a layered decision procedure
for equational implication over magmas that integrates the sound components analyzed
in Section~\ref{sec:main}.

\subsection{Algorithm Overview}

\textsc{predictor4} processes a pair $(E, E')$ of equational laws by a sequence of
tests of increasing computational cost.  Each test either returns a confident
classification or defers to the next layer.  The procedure is summarized as
the following procedure.

\begin{enumerate}[(1)]
  \item \textbf{Trivial checks.}
        If $E'$ is trivial ($\ell_{E'} = r_{E'}$ as terms), or if $E'$ is a
        substitution instance of $E$ (direct specialization), classify as implication.
        If $E$ is trivial, no non-trivial magma satisfies $E$, so the implication
        vacuously fails; classify as non-implication.

  \item \textbf{Knuth-Bendix completion (positive path, sound).}
        Run KB completion on $E$ with the soundness guard
        $\operatorname{vars}(r) \subseteq \operatorname{vars}(\ell)$ and a timeout of $0.8$~seconds.
        If KB terminates with convergent TRS $R$ and
        $\mathrm{NF}_R(\ell_{E'}) = \mathrm{NF}_R(r_{E'})$, classify as implication.
        By Theorem~\ref{thm:kb-soundness}, this classification has no false positives.

  \item \textbf{BFS reachability (positive path, fallback).}
        Expand the set of terms reachable from $\ell_{E'}$ and $r_{E'}$ under
        rewrites by $E$ via breadth-first search, bounded by a step limit.  If the
        reachable sets intersect, classify as implication.  This step is used when
        KB times out or fails; it is faster but provides a weaker soundness guarantee.

  \item \textbf{Linear polynomial counterexample search.}
        Evaluate $E$ and $E'$ in commutative polynomial rings $\mathbb{Z}/p\mathbb{Z}$
        for $15$ primes $p$.  If some evaluation satisfies $E$ but violates $E'$,
        a counterexample is found; classify as non-implication.

  \item \textbf{Random polynomial evaluation.}
        Evaluate $E$ and $E'$ under $640$ random multilinear maps over large primes.
        If any evaluation satisfies $E$ but not $E'$, classify as non-implication.

  \item \textbf{Exhaustive finite enumeration (sizes 2 and 3).}
        Test all $16$ magmas of size 2 and all $19{,}683$ magmas of size 3.
        Additionally test named structures (symmetric group $S_3$, quaternion group
        $Q_8$, cyclic groups, affine magmas).  If any satisfies $E$ but not $E'$,
        classify as non-implication.

  \item \textbf{Z3 SMT model finding (negative path, sound and complete).}
        Solve $\Phi_{E,E',4}$ with a timeout of $0.8$~seconds, then $\Phi_{E,E',5}$
        with $0.6$~seconds.  By Theorem~\ref{thm:z3-complete}, if Z3 returns
        \textsc{sat}, the extracted table $T$ is a verified counterexample; classify
        as non-implication.  If Z3 returns \textsc{unsat}, no size-$n$ counterexample
        exists.

  \item \textbf{Structural feature estimation (fallback).}
        If no test has returned a definitive answer, estimate implication probability
        from structural features (variable counts, term depths, law type, canonical
        signatures).  Return the estimated probability.
\end{enumerate}

\begin{theorem}[No False Positives from KB Component]\label{thm:no-fp}
Let $(E, E')$ be any input pair.  If \textsc{predictor4} classifies $(E, E')$ as an
implication via step~(2) (Knuth-Bendix completion), then $E \models E'$.
\end{theorem}

\begin{proof}
Step~(2) returns a positive classification only if KB completion terminates
successfully with a convergent TRS $R$ satisfying the soundness guard
$\operatorname{vars}(r) \subseteq \operatorname{vars}(\ell)$ for all rules, and the normal-form test
$\mathrm{NF}_R(\ell_{E'}) = \mathrm{NF}_R(r_{E'})$ holds.  By
Theorem~\ref{thm:kb-soundness}, this implies $E \models E'$.
\end{proof}

\begin{theorem}[No False Positives from Z3 Component]\label{thm:z3-no-fp}
If \textsc{predictor4} classifies $(E, E')$ as a non-implication via step~(7),
then $E \not\models E'$.
\end{theorem}

\begin{proof}
Step~(7) returns a negative classification only if Z3 returns \textsc{sat} for
$\Phi_{E,E',n}$ with a concrete satisfying assignment $T$.  This $T$ is a size-$n$
multiplication table for which $E$ holds universally and $E'$ fails existentially;
that is, the magma $(M_T, *)$ with $|M_T| = n$ satisfies $E$ but not $E'$.
Therefore $E \not\models E'$.
\end{proof}

\subsection{Implementation Details}

\para{KB completion.}
The KB completion module is implemented from scratch in Python (version 3.12),
as no Python package providing this algorithm exists on standard package indices.
Key components include: first-order unification via Robinson's algorithm
\cite{robinson1965machine} with occurs check; the Lexicographic Path Order for
rule orientation, with variables ordered alphabetically; critical pair computation
with variable renaming to ensure disjointness; and normal-form reduction with a
term size guard (80 nodes) and step limit (200 steps) to prevent divergence.
The soundness guard ($\operatorname{vars}(r) \subseteq \operatorname{vars}(\ell)$) is checked before each
rule is admitted to $R$.

\para{Z3 encoding.}
The Z3 encoding of $\Phi_{E,E',n}$ uses integer variables $T_{ij}$ constrained
to $\{0,\ldots,n-1\}$.  Term evaluation is encoded recursively: for a subterm
$u * v$, the value under assignment $\varphi$ is
$\mathtt{If}(u_\varphi = 0 \wedge v_\varphi = 0,\, T_{00},\,
   \mathtt{If}(u_\varphi = 0 \wedge v_\varphi = 1,\, T_{01},\, \ldots))$,
giving a linear conditional chain of length $n^2$.  The universally quantified
constraint for $E$ and the existentially quantified violation of $E'$ are expanded
into conjunctions and disjunctions over all assignments.

\para{Environment.}
All experiments use Python~3.12.8, z3-solver~4.16.0, sympy~1.14.0, and
pandas~3.0.2 on an x86-64 CPU.  All random seeds are fixed at 42.
